\section{Tìm hiểu về Process}

\subsection{Khái niệm process trong xv6}

Process là \emph{đơn vị cách ly} của xv6: một process không
thể đọc hay phá bộ nhớ, CPU, file descriptor của process khác, cũng không thể
phá kernel. Để làm được điều này, kernel tạo cho mỗi chương trình hai ảo giác:

\begin{itemize}
  \item \textbf{Bộ nhớ riêng (address space):} mỗi process có một page table
        riêng (\cmd{p->pagetable}) ánh xạ địa chỉ ảo bắt đầu từ 0 sang bộ nhớ
        vật lý. Process này không nhìn thấy bộ nhớ của process khác.
  \item \textbf{CPU riêng (thread):} kernel luân phiên CPU giữa các process.
        Khi một process không chạy, các thanh ghi của nó được lưu lại và khôi
        phục khi nó được chạy tiếp.
\end{itemize}

Trong xv6 mỗi process có đúng một address space và một thread. Mỗi process có
\textbf{hai stack}: user stack dùng khi chạy code user và kernel stack
(\cmd{p->kstack}, 1 trang) dùng khi vào kernel qua system call hoặc
interrupt. Kernel stack được tách riêng và được bảo vệ để kernel vẫn chạy
đúng kể cả khi chương trình user làm hỏng user stack của nó.

Kernel quản lý process bằng một mảng cố định \cmd{struct proc proc[NPROC]}
(\cmd{NPROC = 64}). Không có cấp phát động: tạo process nghĩa là tìm một ô
có \cmd{state == UNUSED} và điền thông tin vào (\cmd{allocproc()}).

\subsection{Các field \texttt{pid}, \texttt{state}, \texttt{context} của \texttt{struct proc}}
\label{sec:proc-fields}

\cmd{struct proc} (\cmd{kernel/proc.h}) chứa toàn bộ trạng thái kernel của
một process. Các field được chia nhóm theo khóa bảo vệ chúng
(Bảng~\ref{tab:proc-groups}).

\begin{table}[H]
\centering
\small
\begin{tabularx}{\textwidth}{>{\raggedright\arraybackslash}p{3.3cm} >{\raggedright\arraybackslash}p{4.3cm} X}
\toprule
\textbf{Khóa bảo vệ} & \textbf{Field} & \textbf{Lý do} \\
\midrule
\cmd{p->lock} & \cmd{state}, \cmd{pid}, \cmd{chan}, \cmd{killed}, \cmd{xstate}
  & Được process khác và scheduler trên CPU khác đọc/ghi. \\
\cmd{wait\_lock} (toàn cục) & \cmd{parent}
  & Process chưa biết cha mình là ai cho tới khi lấy được khóa, nên không thể
    dùng khóa riêng của cha. \\
Không cần khóa & \cmd{kstack}, \cmd{sz}, \cmd{pagetable}, \cmd{trapframe},
  \cmd{context}, \cmd{ofile}, \cmd{cwd}, \cmd{name}
  & Chỉ chính process đó sử dụng. \\
\bottomrule
\end{tabularx}
\caption{Các nhóm field của \texttt{struct proc}}
\label{tab:proc-groups}
\end{table}

\subsubsection{\texttt{pid} -- định danh process}

\cmd{pid} là số nguyên dương duy nhất, cấp bởi \cmd{allocpid()} khi
\cmd{allocproc()} tìm được ô trống. Hàm này tăng biến toàn cục \cmd{nextpid}
dưới khóa \cmd{pid\_lock}, nên hai CPU tạo process cùng lúc không bao giờ nhận
trùng pid. Process đầu tiên (\cmd{init}) có pid 1. pid không được tái sử dụng
dù ô \cmd{proc[]} có được dùng lại.

pid được dùng làm giá trị trả về của \cmd{fork()} cho process cha, giá trị trả
về của \cmd{wait()}, và tham số của \cmd{kill(pid)}. Vì \cmd{kill} và
\cmd{wait} đọc pid của process khác nên field này phải được bảo vệ bởi
\cmd{p->lock}: \cmd{kill} cần chắc chắn process nạn nhân không thoát và bị
cấp phát lại giữa lúc kiểm tra pid và lúc đặt cờ \cmd{killed}.

\subsubsection{\texttt{state} -- trạng thái process}

\cmd{state} có kiểu
\cmd{enum procstate \{UNUSED, USED, SLEEPING, RUNNABLE, RUNNING, ZOMBIE\}},
cho biết ô trong bảng đang trống, đang được dựng, đang chờ sự kiện, sẵn sàng
chạy, đang chạy trên một CPU, hay đã kết thúc nhưng chưa được cha thu dọn.
Hình~\ref{fig:state} và Bảng~\ref{tab:state} mô tả các chuyển trạng thái.

\begin{figure}[H]
\centering
\begin{tikzpicture}[
  st/.style={draw, rounded corners, minimum width=2.3cm, minimum height=0.9cm,
             font=\small\ttfamily, fill=blue!6},
  ar/.style={-{Stealth[length=2.5mm]}, thick},
  lb/.style={font=\scriptsize, fill=white, inner sep=1pt},
  node distance=1.3cm and 2.2cm]
  \node[st] (unused) {UNUSED};
  \node[st, right=of unused] (used) {USED};
  \node[st, right=of used] (runnable) {RUNNABLE};
  \node[st, below=1.9cm of runnable] (running) {RUNNING};
  \node[st, below=1.9cm of used] (sleeping) {SLEEPING};
  \node[st, below=1.9cm of unused] (zombie) {ZOMBIE};

  \draw[ar] (unused) -- node[lb, above]{allocproc} (used);
  \draw[ar] (used) -- node[lb, above]{fork/userinit} (runnable);
  \draw[ar] ([xshift=-4mm]runnable.south) -- node[lb, left]{scheduler} ([xshift=-4mm]running.north);
  \draw[ar] ([xshift=4mm]running.north) -- node[lb, right]{yield} ([xshift=4mm]runnable.south);
  \draw[ar] (running) -- node[lb, below]{sleep} (sleeping);
  \draw[ar] (sleeping) -- node[lb, above, sloped]{wakeup} (runnable);
  \draw[ar] (running.south) -- ++(0,-0.6) -| node[lb, pos=0.25, below]{exit} (zombie.south);
  \draw[ar] (zombie) -- node[lb, left]{wait (cha) + freeproc} (unused);
\end{tikzpicture}
\caption{Sơ đồ chuyển trạng thái của process trong xv6}
\label{fig:state}
\end{figure}

\begin{table}[H]
\centering
\small
\begin{tabularx}{\textwidth}{>{\ttfamily}l X}
\toprule
\textrm{\textbf{Chuyển}} & \textbf{Nguyên nhân} \\
\midrule
UNUSED $\to$ USED & \cmd{allocproc()} tìm thấy ô trống: gán pid, cấp trang
  trapframe, tạo page table, đặt \cmd{context.ra = forkret}. \\
USED $\to$ RUNNABLE & Cuối \cmd{fork()} (hoặc \cmd{userinit()}), sau khi đã
  dựng xong toàn bộ process con. \\
RUNNABLE $\to$ RUNNING & \cmd{scheduler()} chọn process để chạy. \\
RUNNING $\to$ RUNNABLE & \cmd{yield()}: hết lượt do timer interrupt. \\
RUNNING $\to$ SLEEPING & \cmd{sleep(chan)}: chờ sự kiện (đọc pipe, \cmd{wait},
  đọc console, \ldots); \cmd{p->chan} ghi lại kênh đang chờ. \\
SLEEPING $\to$ RUNNABLE & \cmd{wakeup(chan)} từ process khác hoặc interrupt. \\
RUNNING $\to$ ZOMBIE & \cmd{exit()}: đóng file, giao con cho \cmd{init},
  đánh thức cha, lưu \cmd{xstate}. \\
ZOMBIE $\to$ UNUSED & \cmd{wait()} của cha lấy mã thoát rồi gọi
  \cmd{freeproc()}. \\
\bottomrule
\end{tabularx}
\caption{Các chuyển trạng thái và hàm gây ra chúng}
\label{tab:state}
\end{table}

Trạng thái \cmd{ZOMBIE} cần thiết vì hai lý do: process cha
có thể chưa gọi \cmd{wait} khi con thoát nên kernel phải giữ lại mã thoát;
và một process không thể tự giải phóng kernel stack mà nó đang chạy trên đó.
Nếu cha thoát trước con, \cmd{exit()} gọi \cmd{reparent()} để giao các con
cho \cmd{init}, process luôn gọi \cmd{wait} trong vòng lặp, nên mọi process
con đều có người thu dọn.

Trạng thái \cmd{USED} là trạng thái trung gian khi process đang được dựng.
Scheduler chỉ chạy process \cmd{RUNNABLE}, nên một process đang dựng dở bị
``che'' khỏi scheduler cho tới khi \cmd{fork()} đặt nó thành \cmd{RUNNABLE}
ở bước cuối cùng.

\subsubsection{\texttt{context} -- ngữ cảnh kernel thread}

\cmd{struct context} lưu 14 thanh ghi của \emph{kernel thread} của process:
\cmd{ra}, \cmd{sp} và các thanh ghi callee-saved \cmd{s0}--\cmd{s11}. Nó
được dùng bởi hàm assembly \cmd{swtch(old, new)} (\cmd{kernel/swtch.S}):
lưu các thanh ghi hiện tại vào \cmd{old}, nạp các thanh ghi từ \cmd{new},
rồi \cmd{ret}.

\begin{itemize}
  \item \cmd{swtch} chỉ lưu các thanh ghi callee-saved, vì trình biên dịch C đã
        lưu các thanh ghi caller-saved lên stack trước khi gọi hàm.
  \item \cmd{swtch} không lưu \cmd{pc}, mà lưu \cmd{ra}. Khi \cmd{ret}, CPU nhảy tới \cmd{ra} vừa
        nạp, tức là ngay sau lệnh gọi \cmd{swtch} trước đây của thread kia.
  \item Vì \cmd{sp} cũng được nạp lại, CPU chuyển sang chạy trên stack của
        thread kia.
\end{itemize}

Mỗi CPU cũng có một \cmd{context} trong \cmd{struct cpu}, là ngữ cảnh của
\emph{scheduler thread} của CPU đó. Process chuyển CPU bằng
\cmd{swtch(\&p->context, \&mycpu()->context)} và scheduler chuyển CPU cho
process bằng \cmd{swtch(\&c->context, \&p->context)}.

Cần phân biệt \cmd{context} với \cmd{trapframe} (Bảng~\ref{tab:ctx-tf}):

\begin{table}[H]
\centering
\small
\begin{tabularx}{\textwidth}{l X X}
\toprule
 & \cmd{p->trapframe} & \cmd{p->context} \\
\midrule
Lưu gì & Toàn bộ thanh ghi \emph{user} (\cmd{epc}, \cmd{sp}, \cmd{a0}--\cmd{a7}, \ldots)
       & \cmd{ra}, \cmd{sp}, \cmd{s0}--\cmd{s11} của \emph{kernel thread} \\
Khi nào & User $\to$ kernel (system call, interrupt)
        & Kernel thread $\to$ scheduler (\cmd{swtch}) \\
Ai lưu & \cmd{uservec} trong \cmd{trampoline.S} & \cmd{swtch.S} \\
Mục đích & Quay về đúng lệnh trong chương trình user
         & Quay về đúng chỗ trong code kernel \\
\bottomrule
\end{tabularx}
\caption{So sánh \texttt{trapframe} và \texttt{context}}
\label{tab:ctx-tf}
\end{table}

Khi \cmd{allocproc()} tạo process mới, nó đặt \cmd{context.ra = forkret} và
\cmd{context.sp = p->kstack + PGSIZE}. Do đó lần đầu scheduler \cmd{swtch} vào
process mới, CPU ``quay về'' đầu hàm \cmd{forkret()} trên kernel stack rỗng
của process đó.

\subsection{Trace system call \texttt{fork}}
\label{sec:trace-fork}

\subsubsection{Từ user vào kernel}

\begin{enumerate}
  \item Chương trình user gọi \cmd{fork()}. Hàm này được sinh bởi
        \cmd{user/usys.pl}: nạp số hiệu \cmd{SYS\_fork} vào thanh ghi \cmd{a7}
        rồi thực thi lệnh \cmd{ecall}.
  \item \cmd{ecall} chuyển CPU sang supervisor mode và nhảy tới \cmd{uservec}
        (\cmd{kernel/trampoline.S}). \cmd{uservec} lưu toàn bộ thanh ghi user
        vào \cmd{p->trapframe}, chuyển sang page table kernel và kernel stack
        của process, rồi gọi \cmd{usertrap()}.
  \item \cmd{usertrap()} (\cmd{kernel/trap.c}) thấy \cmd{scause == 8}
        (system call), tăng \cmd{trapframe->epc} thêm 4 để khi trở về sẽ chạy
        lệnh sau \cmd{ecall}, bật interrupt và gọi \cmd{syscall()}.
  \item \cmd{syscall()} đọc \cmd{a7}, tra bảng \cmd{syscalls[]} và gọi
        \cmd{sys\_fork()}, hàm này chỉ gọi \cmd{fork()} trong \cmd{kernel/proc.c}.
        Giá trị trả về được ghi vào \cmd{trapframe->a0}.
\end{enumerate}

\subsubsection{Hàm \texttt{fork()} trong kernel}

\begin{enumerate}
  \item \textbf{Cấp phát process:} \cmd{allocproc()} duyệt \cmd{proc[]}, tìm
        ô \cmd{UNUSED}, gán pid mới, chuyển sang \cmd{USED}, cấp trang
        trapframe, tạo page table rỗng (chỉ có trampoline và trapframe) và
        khởi tạo \cmd{context}. Hàm trả về với \cmd{np->lock} đang được giữ.
        Nếu hết ô hoặc hết bộ nhớ, \cmd{fork} trả về $-1$.
  \item \textbf{Sao chép bộ nhớ:} \cmd{uvmcopy()} cấp trang mới và chép toàn
        bộ bộ nhớ user (code, data, stack, heap) của cha sang con, rồi đặt
        \cmd{np->sz = p->sz}. Hai process có nội dung bộ nhớ giống nhau nhưng
        nằm ở các trang vật lý khác nhau.
  \item \textbf{Sao chép thanh ghi user:} \cmd{*(np->trapframe) = *(p->trapframe)}.
        Do đó khi con trở về user space, nó tiếp tục chạy ngay sau lệnh
        \cmd{ecall} của \cmd{fork}, giống hệt cha.
  \item \textbf{Giá trị trả về của con:} \cmd{np->trapframe->a0 = 0}. Đây là
        điểm khác duy nhất giữa thanh ghi của cha và con, giải thích vì sao
        \cmd{fork} ``trả về hai lần'': 0 ở con, pid của con ở cha.
  \item \textbf{Chia sẻ file đang mở:} với mỗi \cmd{p->ofile[i]} khác 0,
        gán \cmd{np->ofile[i]} bằng \cmd{filedup(p->ofile[i])}. \cmd{filedup} chỉ tăng
        \cmd{ref} của \cmd{struct file} chứ không mở file mới, nên cha và con
        dùng chung offset. Thư mục hiện tại được chia sẻ bằng \cmd{idup(p->cwd)}.
  \item \textbf{Tên:} chép \cmd{p->name} sang \cmd{np->name}.
  \item \textbf{Ghi nhận cha:} nhả \cmd{np->lock}, lấy \cmd{wait\_lock}, đặt
        \cmd{np->parent = p}.
  \item \textbf{Cho phép chạy:} lấy lại \cmd{np->lock}, đặt
        \cmd{np->state = RUNNABLE}. Chỉ từ lúc này scheduler mới có thể chọn
        process con.
  \item Trả về pid của con. Giá trị này đi qua \cmd{trapframe->a0} của cha.
\end{enumerate}

\subsubsection{Lần chạy đầu tiên của process con}

Khi một CPU chọn process con, \cmd{scheduler()} gọi
\cmd{swtch(\&c->context, \&np->context)}. Vì \cmd{context.ra = forkret}, CPU
nhảy vào \cmd{forkret()}. Hàm này nhả \cmd{p->lock} (khóa được scheduler lấy
trước khi \cmd{swtch}, xem mục~\ref{sec:lock-swtch}), lần đầu tiên thì khởi
tạo file system, rồi gọi \cmd{usertrapret()}. \cmd{usertrapret()} và
\cmd{userret} nạp lại thanh ghi user từ trapframe đã sao chép ở bước 3 và
thực thi \cmd{sret}. Process con trở về user space ngay sau lệnh \cmd{ecall},
với \cmd{a0 = 0}.

Process cha trở về theo đường bình thường: \cmd{syscall()} ghi pid vào
\cmd{a0}, \cmd{usertrap()} gọi \cmd{usertrapret()} và cha tiếp tục chạy. Thứ
tự cha hay con chạy trước phụ thuộc vào scheduler.

\subsection{Cơ chế lập lịch: \texttt{scheduler()}, \texttt{sched()}, \texttt{yield()}}
\label{sec:sched}

\subsubsection{Khi nào CPU chuyển sang process khác}

xv6 chuyển process trong hai trường hợp:
\begin{itemize}
  \item \textbf{Tự nguyện:} process gọi một system call phải chờ (\cmd{read},
        \cmd{wait}, \cmd{sleep}, \ldots) và đi vào \cmd{sleep()}.
  \item \textbf{Bắt buộc:} timer interrupt định kỳ. Khi \cmd{devintr()} trả
        về 2 (timer), \cmd{usertrap()} hoặc \cmd{kerneltrap()} gọi
        \cmd{yield()}, ép process đang tính toán nhường CPU.
\end{itemize}

\subsubsection{\texttt{scheduler()}}

Mỗi CPU, sau khi khởi động xong, gọi \cmd{scheduler()} và không bao giờ
quay lại. Hàm này là một vòng lặp vô hạn chạy trên stack riêng của CPU
(scheduler thread):

\begin{enumerate}
  \item Bật interrupt (tránh deadlock khi mọi process đều đang chờ I/O).
  \item Duyệt lần lượt \cmd{proc[0..NPROC-1]}. Với mỗi \cmd{p}: lấy
        \cmd{p->lock}; nếu \cmd{p->state == RUNNABLE} thì đặt
        \cmd{p->state = RUNNING}, \cmd{c->proc = p}, rồi
        \cmd{swtch(\&c->context, \&p->context)} để chạy \cmd{p}.
  \item Khi \cmd{p} nhường CPU, \cmd{swtch} ``trả về'' trong scheduler. Đặt
        \cmd{c->proc = 0}, nhả \cmd{p->lock} và xét process tiếp theo.
  \item Nếu duyệt hết bảng mà không có process nào \cmd{RUNNABLE}, CPU thực thi
        \cmd{wfi} để nghỉ tới interrupt tiếp theo.
\end{enumerate}

Chính sách này là \textbf{round robin}: các process được chạy lần lượt,
không có độ ưu tiên. Scheduler chạy trên stack riêng thay vì
stack của một process, vì nếu dùng stack của process thì process đó có thể
được đánh thức và chạy trên CPU khác, dẫn tới hai CPU dùng chung một stack.

\subsubsection{\texttt{sched()}}

\cmd{sched()} là \emph{điểm duy nhất} một kernel thread nhường CPU cho
scheduler. Trước khi gọi \cmd{swtch(\&p->context, \&mycpu()->context)}, nó
kiểm tra bốn điều kiện, vi phạm bất kỳ điều kiện nào sẽ gây \cmd{panic}:
\begin{itemize}
  \item đang giữ \cmd{p->lock};
  \item chỉ giữ đúng một khóa (\cmd{mycpu()->noff == 1}), vì giữ thêm khóa
        khác qua \cmd{swtch} có thể gây deadlock;
  \item \cmd{p->state != RUNNING}, tức là người gọi đã đổi trạng thái;
  \item interrupt đang tắt.
\end{itemize}
\cmd{sched()} cũng lưu và khôi phục \cmd{intena}, vì đây là thuộc tính của
kernel thread chứ không phải của CPU. Mọi nơi muốn nhường CPU (\cmd{yield},
\cmd{sleep}, \cmd{exit}) đều làm theo cùng một mẫu: lấy \cmd{p->lock}, đổi
\cmd{p->state}, gọi \cmd{sched()}.

\subsubsection{\texttt{yield()}}

\cmd{yield()} là cách nhường CPU khi hết lượt: lấy \cmd{p->lock}, đặt
\cmd{p->state = RUNNABLE}, gọi \cmd{sched()}. Khi process được chọn chạy lại,
\cmd{sched()} trả về và \cmd{yield()} nhả \cmd{p->lock}. Process ở trạng
thái \cmd{RUNNABLE} nên sẽ được scheduler chọn lại ở một vòng duyệt sau.

\subsubsection{Toàn bộ quá trình chuyển process}

Hình~\ref{fig:switch} mô tả việc
chuyển từ process \cmd{shell} sang process \cmd{cat} trên một CPU. Mỗi lần
chuyển gồm \textbf{hai} lần \cmd{swtch}: process $\to$ scheduler và
scheduler $\to$ process mới. \cmd{sched} và \cmd{scheduler} luân phiên
chuyển quyền điều khiển cho nhau, nên chúng được gọi là các \emph{coroutine}.

\begin{figure}[H]
\centering
\begin{tikzpicture}[
  box/.style={draw, rounded corners, font=\scriptsize\ttfamily, align=center,
              minimum height=0.75cm, fill=blue!5},
  ar/.style={-{Stealth[length=2mm]}, thick},
  lb/.style={font=\scriptsize, align=center}]
  % user / kernel boundary
  \draw[dashed, gray] (-1.2,0) -- (13.2,0);
  \node[font=\scriptsize, gray] at (6,0.25) {user space};
  \node[font=\scriptsize, gray] at (6,-0.25) {kernel space};

  \node[box, fill=green!8] (sh) at (0,1.1) {shell\\(user)};
  \node[box] (ut) at (0,-1.2) {usertrap};
  \node[box] (y) at (0,-2.5) {yield\\sched};
  \node[box, fill=orange!12] (sc) at (6,-2.5) {scheduler\\(stack CPU)};
  \node[box] (s2) at (12,-2.5) {sched\\yield};
  \node[box] (ur) at (12,-1.2) {usertrapret};
  \node[box, fill=green!8] (cat) at (12,1.1) {cat\\(user)};

  \draw[ar] (sh) -- node[lb, right, pos=0.25, fill=white]{timer interrupt} (ut);
  \draw[ar] (ut) -- (y);
  \draw[ar] (y) -- node[lb, above]{swtch(\&shell->context,\\\&cpu->context)} (sc);
  \draw[ar] (sc) -- node[lb, above]{swtch(\&cpu->context,\\\&cat->context)} (s2);
  \draw[ar] (s2) -- (ur);
  \draw[ar] (ur) -- node[lb, left, pos=0.75, fill=white]{sret} (cat);

  \node[lb, below=0.15cm of y] {kstack shell};
  \node[lb, below=0.15cm of sc] {release(shell->lock)\\acquire(cat->lock)};
  \node[lb, below=0.15cm of s2] {kstack cat};
\end{tikzpicture}
\caption{Chuyển CPU từ \texttt{shell} sang \texttt{cat}}
\label{fig:switch}
\end{figure}

\subsubsection{Khóa \texttt{p->lock} được giữ qua \texttt{swtch}}
\label{sec:lock-swtch}

Thông thường thread nào lấy khóa thì chính thread đó nhả khóa. Trong xv6,
\cmd{p->lock} được giữ xuyên qua \cmd{swtch}:
\begin{itemize}
  \item \cmd{yield()} của \cmd{shell} lấy \cmd{shell->lock}, nhưng
        \cmd{scheduler()} mới là nơi nhả nó.
  \item \cmd{scheduler()} lấy \cmd{cat->lock}, nhưng \cmd{yield()} của
        \cmd{cat} (hoặc \cmd{forkret()} nếu \cmd{cat} vừa được tạo) mới nhả nó.
\end{itemize}
Lý do: nếu \cmd{shell->lock} được nhả ngay sau khi đặt
\cmd{state = RUNNABLE} nhưng trước khi \cmd{swtch} lưu xong ngữ cảnh, scheduler
của một CPU khác có thể chọn \cmd{shell} và chạy nó trên chính kernel stack
mà CPU hiện tại vẫn đang dùng. Giữ khóa tới khi \cmd{swtch} hoàn tất đảm bảo
mỗi process chỉ chạy trên đúng một CPU tại một thời điểm.

\subsubsection{Liên hệ với phần cài đặt \texttt{sysinfo}}

Để trả về ``number of scheduling events'' theo yêu cầu 4.2, nhóm thêm field
\cmd{nsched} vào \cmd{struct proc}:
\begin{itemize}
  \item Khởi tạo bằng 0 trong \cmd{allocproc()}, vì ô \cmd{proc[]} được dùng
        lại cho process mới.
  \item Tăng 1 trong \cmd{scheduler()} ngay sau \cmd{p->state = RUNNING} và
        trước \cmd{swtch}. Đây là nơi duy nhất một process được cấp CPU, nên
        \cmd{nsched} đếm chính xác số lần process được lập lịch. Tại điểm này
        scheduler đang giữ \cmd{p->lock}, nên hai CPU không thể cùng tăng biến.
\end{itemize}
Khi đọc thông tin process, \cmd{sysinfo} lấy \cmd{wait\_lock} trước rồi mới
lấy \cmd{p->lock}, cùng thứ tự với \cmd{wait()} và \cmd{exit()}, để đọc
\cmd{parent} đúng mà không gây deadlock.
